\documentclass[../main.tex]{subfiles}
\begin{document}
%==========================================TIÊU ĐỀ TẬP========================================
\begin{table}[h]
  \begin{adjustwidth}{-1cm}{}
    \begin{tabular}{>{\centering\arraybackslash}p{6cm}|>{\raggedright\arraybackslash}p{12.5cm}}
      \multirow{5}{6cm}
      % Cột bên trái: ÉPISODE và số tập
      {\\[-40pt]
        \flaregothic\fontsize{20pt}{20pt}\selectfont \centering
        \color{eptype}ÉPISODE\color{black}\\
        \burbank\fontsize{72pt}{72pt}\selectfont
        \color{epnum}60\color{black}\\[6pt]
        \cafeteria\fontsize{20pt}{20pt}\selectfont\color{epnum}(5-4)\color{black}
        \\[18pt]
        \flaregothic\fontsize{18pt}{18pt}\selectfont
        \color{epattr}ÉPISODE PERDU
        \color{black}
        \\}
       &
      % Dòng thứ nhất: tên tập tiếng Nhật
      {\fotskip\fontsize{18pt}{18pt}\selectfont とどめの\ruby{一}{いち}\ruby{撃}{げき}を\ruby{刺}{さ}したのは\ruby{誰}{だれ}だ！？\ruby{罪}{つみ}のなすりつけ\ruby{合}{あ}い！
      } \\[6pt]
      % Dòng thứ hai: tên tập tiếng Anh
       & {\cafeteria\fontsize{18pt}{18pt}\selectfont \hll\ Closed-Room Mystery} \\[4pt]
      % Dòng thứ ba: tên tập tiếng Việt
       & {\vnmsans\fontsize{18pt}{18pt}\selectfont Đổ tội cho nhau} \\[8pt]
      % Dòng thứ tư: Nhân vật xuất hiện (theo đúng thứ tự debut: OC → Roboco(11) → Haato(19) → Shion(21) → Korone(27))
       & {\caviar{\fontsize{14pt}{14pt}\selectfont Nhân vật: \emp{kuon1} \emp{roboco} \emp{haato} \emp{shion} \emp{korone}}}
      \\[8pt]
      % Dòng cuối cùng: Ngày phát hành
       & {\continuum\fontsize{16pt}{16pt}\selectfont Ngày phát hành: 28/06/2020} \\[18pt]
    \end{tabular}
  \end{adjustwidth}
\end{table}
%=========================================TRUYỆN CHÍNH========================================
\color{black}

% Hộp tóm tắt
\txtbx[2]{{\color{vol} \uline{Tóm tắt:}} Một buổi sáng mùa hè nóng nực, Shion và Haato phát hiện thứ trông giống Roboco đang nằm bất động trên sàn, rỉ ra chất lỏng màu vàng cam. Hoảng loạn, họ lau dọn, đổ dầu ô liu vào, rồi cãi vã đổ tội cho nhau, kéo theo cả Korone và Kuon vào cuộc chiến. Khi Roboco thật trở về, họ mới vỡ lẽ: đó chỉ là bình đựng cà ri mới của cô. Họ ném bình ra cửa sổ, làm vỡ kính, rồi cuối cùng cùng nhau thưởng thức cà ri, mà không ai kể cho Roboco biết về vụ dầu ô liu.}

Cuối tháng Sáu, tiết trời oi bức như một chiếc lò nướng khổng lồ. Độ ẩm cao đến mức không khí như đặc quánh lại, chỉ cần đưa tay lên là có thể cảm nhận được hơi nước đang bám vào da. Những tiếng ve kêu râm ran từ bên ngoài vọng vào, như một dàn đồng ca bất tận, nhắc nhở mọi người rằng mùa hè đã thực sự đến, và nó chẳng hề dễ chịu chút nào.

Bên trong văn phòng \hll, cái nóng cũng chẳng kém cạnh. Những cái quạt cây dù đã chạy hết công suất nhưng vẫn chỉ đủ để khiến mọi người không bị tan chảy. Kuon đã có mặt từ sớm, tranh thủ ngồi nghỉ một lát sau khi xử lý một đống giấy tờ cao ngất. Mồ hôi vẫn lấm tấm trên trán, nhưng ít nhất cậu đã hoàn thành xong báo cáo, một chiến tích đáng tự hào trong ngày hè nóng nực.

Không lâu sau, cánh cửa bật mở. Shion và Haato bước vào. Hai người chào buổi sáng cậu bằng giọng đều đều, cũng mệt mỏi như chính thời tiết. Cậu cũng đáp lại họ như mọi ngày.

``Chào em. Hôm nay cũng nóng nhỉ.''

``Nóng chết đi được\dots''  Shion vừa nói vừa quạt tay, bước vào trong.

Haato thì lăn ra ghế, mặt úp xuống bàn: ``Em muốn ở nhà tắm mát\dots\ sao phải đến công ty chứ\dots''

Bỗng nhiên, nàng Phù thủy phát hiện ra một thứ gì đó kỳ lạ ngay dưới sàn, gần chân bàn làm việc của Kuon. Cô khẽ ``hửm'' một tiếng, rồi cúi xuống quan sát kỹ hơn. Đôi mắt cô nheo lại, vừa tò mò vừa hoang mang.

Đó là Roboco - hoặc ít nhất cô nghĩ vậy - đang ngồi với tư thế kỳ lạ: hai tay ôm mặt, đầu gập xuống, ngồi co ro như một đứa trẻ bị la mắng. Không một cử động, không một âm thanh. Trông cô như một bức tượng sáp trong bảo tàng, nếu không thì là\dots\ một xác ướp mới được khai quật.

\img{5-4a}

``Roboco-san bị sao vậy?''  Shion chỉ vào người mà cô cho là nàng Người máy, giọng đầy lo lắng.

Cậu Tổng nhún vai, mặt vẫn bình thản: ``Chịu. Vốn dĩ đã như thế này khi anh bước vào rồi. Anh còn tưởng em ấy đang thực hành thiền định hoặc kiểu\dots\ sạc pin bằng tư thế này.''

Nàng Song tính, từ một góc độ khác, lên giả thuyết với vẻ mặt suy tư của một nhà khoa học nghiệp dư: ``Có vẻ như chị ấy bị gỉ sét do độ ẩm vào mùa này, nên đang giữ tư thế thoải mái nhất thôi. Giống như máy móc lâu ngày không dùng, bôi trơn một chút lại chạy ngon.''

Cậu Tổng nhướng mày: ``Cũng có thể chứ. Mà anh nhớ là Robo-ngố này làm bằng kim loại chống ăn mòn mà nhỉ?''

Shion tiếp nhìn với sự nghi ngờ, tay sờ cằm: ``Em cũng nhớ là thế mà\dots\ Sao chị ấy lại không hoạt động thế này? Hay là hết pin?''

Nói rồi, cô tiến tới và chạm nhẹ vào vai nàng Rô-bốt đang nằm im trên sàn. Một cái chạm nhẹ, gần như chỉ là đặt tay lên.

Không ai ngờ được.

Ngay khi vừa chạm vào, toàn bộ cơ thể Roboco - vốn đang ngồi ở tư thế co ro - bất ngờ đổ kềnh ra phía sau, giống như một chiếc bánh domino vừa bị hất đổ. Toàn thân vẫn đông cứng lại, chẳng hề nhích lấy một ly.

% \img{5-4b}

``\textbf{Éc!?}''  Shion thảng thốt, vội lùi ra xa gần hai mét, mặt tái mét.

Haato đang nằm trên ghế cũng bật dậy, sồn sồn nhảy vào, mắt trợn tròn: ``G-Gì vậy? Tiếng gì thế!? Có chuyện gì xảy ra?''

Nàng Phù thủy vội lên tiếng trấn an tiền bối của mình, dù giọng vẫn còn run run: ``B-Bình tĩnh nào! Không có gì đâu! Gió thổi ngã đó mà\dots\ Trời đang nóng, gió lùa vào, vậy thôi.''

Cậu Tổng nhìn cô đầy mỉa mai, khóe môi nhếch lên: ``Gió gì mà mạnh ghê nhỉ. Đủ sức thổi đổ một người máy nặng cả tạ? Chắc phải gió cấp 12.''

Shion liếc cậu, buông ra một câu đầy uy lực: ``Câm miệng đi. Đừng có châm biếm.''

Cậu chỉ nhún vai, giơ tay đầu hàng.

Ngay lúc đó, một thứ chất lỏng màu vàng cam bất ngờ chảy ra từ\dots\ từ đâu đó bên dưới người Roboco. Nó chảy thành một vệt dài trên sàn, loang ra từ từ, tạo nên một hình dạng trông như một bản đồ vô định.

\img{5-4c}

Nàng Song tính hét lên trong hoảng loạn, giọng xé toạc bầu không khí ngột ngạt: ``\textbf{Chất lỏng gì thế!? Cái thứ này chảy ra từ người chị ấy là cái gì vậy!? Máu? Dầu? Hay là\dots\ nước ép cà chua?}''

Shion vội vàng phân trần, tay vẫy vẫy như thể đang xua đuổi một con ruồi: ``C-Chắc là cà ri thôi! Chị ấy đã được tắm trong cà ri trong buổi stream ngày hôm qua, thế là\dots\ có thể còn sót lại trên người, rồi lúc đổ ra\dots\ ừm, đó là lý do hợp lý mà!''

Haachama lập tức phản bác, mắt nheo lại đầy nghi ngờ: ``Em không lừa được ai đâu! Chắc chắn đây là dầu mà! Dầu máy, dầu nhớt, thứ mà người máy cần để bôi trơn! Chị ấy bị rò rỉ dầu rồi! Cần gọi thợ sửa ngay!''

Cậu Tổng cố gắng bào chữa cho Shion, giọng tỉnh bơ nhưng đầy thiện chí: ``Shion-chan nói đúng mà, hôm qua anh xem stream của em ấy, quả thực có một buổi tắm cà ri khá\dots\ hoành tráng. Nước cà ri màu vàng cam, chảy xuống sàn, trông cũng giống thế này.''

Nhưng rồi, người mà cậu đang cố gắng bào chữa lại chẳng quan tâm đến lời cậu. Cô đang hoảng loạn chạy khắp phòng, mắt dáo dác tìm kiếm: ``Khăn giấy đâu? Giẻ lau đâu? Có ai mang vào đây không? Nhanh lên, lau đi, kẻo lem hết sàn nhà, lát Fubuki-senpai vào thấy lại bảo bọn mình phá hoại!''

Kuon chỉ có thể biết đường thở dài, lẩm bẩm một mình: ``Có ai thèm nghe mình nói đâu\dots\ toàn tự biên tự diễn.''

Ngay sau đó, Haato và Shion,  mỗi người cầm một cuộn giấy ăn, đồng thanh quát về phía cậu: ``Cả anh nữa! Đừng có đứng ngây ra đó! Đến phụ lau đi chứ!''

Cậu lại thở dài lần nữa, đứng dậy khỏi ghế, lững thững đi tìm cây chổi lau nhà trong góc bếp. Đằng sau lưng cậu, hai cô gái đang hì hục lau vệt chất lỏng màu vàng cam, thứ mà có thể là cà ri, có thể là dầu, hoặc có thể là một thứ nguyên tố bí ẩn nào đó còn nguy hiểm hơn thế.

\ct

Sau một hồi lau dọn mướt mồ hôi, cả ba người đều nhận ra một sự thật nghiệt ngã: thứ chất lỏng màu vàng cam đó trào ra quá nhiều, đến mức họ không thể nào lau xuể. Những cuộn giấy ăn lần lượt trở nên ướt nhẹp, thấm đẫm, rồi biến thành những cục bột nhão vô dụng. Cây chổi lau nhà cũng chẳng giúp ích được gì, nó chỉ quệt đều chất lỏng ra một vùng rộng hơn, biến sàn nhà thành một tấm bản đồ mờ ảo.

Shion đứng thẳng dậy, lưng mỏi nhừ, nhìn xấp giấy ăn cuối cùng đã nhuốm màu vàng mà bất lực: ``Làm sao bây giờ\dots\ Nó như có mạch nước ngầm ấy, lau mãi không hết.''

Haato cũng không kém cạnh, cô đang quỳ gối, hai tay ôm đầu: ``Chảy ra hết rồi\dots\ Càng lau càng nhiều, cứ như thể Roboco-san đang chảy nước mắt vậy\dots\ mà nước mắt màu cà ri.''

Cả hai cô gái vội vàng lui ra xa, quay mặt vào nhau thì thầm bàn bạc, mắt liếc ngang liếc dọc như những kẻ trộm đang lên kế hoạch che giấu tang vật.

``Chúng ta nên giữ bí mật về chuyện này, đừng để cho ai biết\dots''  Shion thì thầm.

``Chị đồng ý. Nếu bị phát hiện, có khi chúng ta sẽ bị bắt đi lau dọn cả ngày mất! Hôm qua em mới làm đổ nước ngọt ra bàn, A-chan đã la om sòm rồi.''

Trong khi hai cô gái đang bàn kế ``phi tang'', cậu Tổng Kuon tiến lại gần đống chất lỏng đang chảy loang trên sàn. Cậu không giấu nổi sự thắc mắc, cúi xuống quan sát thật kỹ.

``\textit{Trông nó có vẻ không giống dầu lắm\dots\ Dầu thường sóng sánh, trong mờ, còn thứ này sệt hơn, màu sắc cũng đặc biệt\dots}''

Cậu đưa mũi tiến gần vết lỏng, ngửi nhẹ. Một mùi hương quen thuộc, nồng ấm, thoang thoảng gia vị xộc vào mũi. Mắt cậu hơi nheo lại.

``\textit{Sao mùi giống cà ri thế nhỉ? Không lẽ\dots}''

Cậu khẽ cúi xuống, dùng ngón tay phết một ít chất lỏng lên đầu ngón tay, đưa lên ngửi kỹ hơn. Mùi thơm đặc trưng của nghệ, thì là, và các loại gia vị vẫn còn nguyên vẹn. Như để xác nhận lần cuối, cậu liều lĩnh đưa lên miệng nếm thử một chút.

Một thoáng im lặng.

Rồi mắt cậu mở to, long lanh như hai viên bi, khuôn mặt biến sắc từ nghi ngờ sang sửng sốt: ``\textbf{\textit{C-Cái này là cà ri mà!}}''

Cậu liếc nhìn thứ gì đó trông giống như Roboco đang nằm bất động trên sàn, đầu óc nhanh chóng lắp ghép các mảnh ghép. Một giả thuyết kỳ quá  nhưng có lẽ là đúng hiện ra trong đầu: ``\textit{\dots\ Không lẽ\dots\ đây là bình chứa cà-ri của PON-chan thì sao? Tại sao lại\dots\ mà cũng hợp lý thật, hôm qua em ấy tắm cà ri xong chưa kịp vệ sinh, chất lỏng còn sót lại trong các khe máy, lâu ngày bị nóng chảy ra?}''

Cậu chưa kịp hoàn thành chuỗi suy luận thì tiếng cửa mở đột ngột vang lên, cắt ngang cuộc họp khẩn của ba người.

Cánh cửa mở ra, và Korone xông vào. Cô hớn hở cất giọng chào, như thể chẳng có gì bất thường xảy ra: ``Chào buổi sáng\~{} Hôm nay trời nóng ghê, mọi người có khỏe không?''

Chẳng kịp để nàng Chó hiểu đầu cua tai nheo gì đang xảy ra, hai bóng người lao về phía cô với vẻ mặt tuyệt vọng như những người sắp chết đuối vớ được phao. Shion và Haachama sà vào Korone, nước mắt lưng tròng, tay chắp lại cầu xin như thể vừa gặp được cứu tinh.

``Korone-chan! Cứu tụi em với!''  Shion run rẩy.

``Koro-chan, chị chỉ có mình em thôi!''  Haato cũng không kém phần thảm thiết.

\img{5-4e}

Korone tròn mắt, đôi tai chó dựng đứng lên vì ngạc nhiên. Cô nhìn sang cậu con trai đang đứng gần đó - người duy nhất trông vẫn còn bình tĩnh - rồi hỏi: ``Có chuyện gì thế, Kuon-senpai? Sao hai người ấy khóc lóc thảm thiết vậy?''

Cậu Tổng chỉ thở dài, bắt đầu kể lại toàn bộ sự việc từ đầu: chuyện Roboco nằm bất động, chuyện chất lỏng màu vàng cam tràn ra, chuyện cậu vừa nếm thử và xác nhận đó là cà ri (mà ngay lập tức bị hai người kia phản bác). Hai cô gái bên cạnh thi thoảng góp lời, giọng đầy hoảng loạn, mỗi người một kiểu, khiến câu chuyện càng thêm rối rắm.

``Tóm lại là có dầu tràn ra ngoài à?''  Korone gật gù, dường như đã hiểu được sự tình.

Cậu Tổng nghe vậy, xua tay định đính chính: ``Koro-chan, anh không nghĩ đó là dầu đâu, anh vừa nếm thử rồi, nó là cà--''

Nhưng lập tức bị nàng Phù thủy cắt ngang, giọng như cầu cứu: ``Ta phải làm gì bây giờ? Chúng ta có nên giấu xác\dots\ à nhầm, giấu bằng chứng không?''

Korone nghe vậy, tỉnh bơ phán một câu, đôi mắt sáng lên như một vị thẩm phán: ``Hai người phải thú tội thôi.''

Cậu quản lý đồng tình, gật đầu: ``Ừ, anh nghĩ đó cũng là cách tốt nhất đấy. Dù sao Roboco-chan không phải là người máy dễ rỉ dầu, lỡ có chuyện gì thì thà nói thật còn hơn để người khác phát hiện.''

Haato và Shion nhìn nhau, mặt xanh mét. Họ chưa kịp phản ứng thì Korone, với vẻ mặt vô cùng nghiêm túc, lôi từ trong túi áo khoác ra hai món đồ khiến cậu con trai phải giật bắn mình.

Một cây thánh giá bằng gỗ nhỏ, và một chiếc đinh dài.

Cô nở một nụ cười đầy ẩn ý, ánh mắt sắc như dao: ``Đây, một cây thánh giá và đinh. Các chị hãy quỳ xuống, thú nhận tội lỗi, và em sẽ đóng đinh để xua đuổi tà ma theo nghi lễ cổ xưa thôi.''

\img{5-4f}

Haato thảng thốt, lùi lại một bước, mắt mở to: ``Đ-Đóng đinh sao!? Bọn chị có làm gì sai đâu! Có cần phải hành hình không hả!?''

Shion thì lại run rẩy, mắt cô dần mất đi ánh sáng, thẫn thờ nhìn đôi tay đang run lên của mình. Cô thều thào: ``Haato-chan\dots\ Em sẽ đóng đinh vào tay em được chứ? Nếu phải chịu hình phạt, em xin nhận\dots''

% \img{5-4g}

Nhìn phản ứng có phần hơi quá đà của cô, nàng Song tính thốt lên, giọng đầy bất lực: ``Chị không có làm đâu! Shion-chan, em bình tĩnh lại đi! Chưa gì đã muốn tự tử rồi!''

Cậu Tổng chỉ ngán ngẩm thở dài, đưa tay xoa trán: ``Có phải là Chúa Giê-su đâu mà làm thế hả em\dots\ Thú tội là lên tiếng nhận lỗi, chứ không phải lên thập tự giá. Với lại, chúng ta đang ở văn phòng, không phải nhà thờ.''

Korone vẫn giữ nụ cười, nhưng hạ tay xuống, cất thánh giá và đinh vào túi. Cô cười khì: ``Đùa thôi mà. Nhưng các chị nên nhanh lên, lát nữa có người vào thì hơi muộn.''

\ct

Sau một hồi giúp nàng Phù thủy lấy lại tinh thần (và thuyết phục cô rằng mình không phải Chúa Giê-su), Korone rốt cuộc cũng ngồi xuống lắng nghe kế hoạch của Haachama. Nàng Song tính đứng giữa phòng, hai tay chống nạnh, mắt sáng rực như một nhà phát minh vừa tìm ra lời giải cho bài toán vũ trụ: ``Nghe này! Nếu chị ấy mất dầu thì chúng ta cứ đổ đầy lại thôi! Đơn giản, logic, và hiệu quả!''

Cô nói, tràn đầy sự tự tin\dots\ tếu. Có lẽ đây là kế hoạch duy nhất mà cô nghĩ ra được trong lúc hoảng loạn, nhưng ít nhất thì nó cũng là một kế hoạch.

Trước mặt cô, Shion và Korone - dù chẳng muốn chút nào - đang ngồi lắng nghe, mặt mỗi người mỗi vẻ. Nàng Phù thủy thì gật gù như một học sinh ngoan ngoãn, còn nàng Khuyển thì trông như đang nhai kẹo cao su và nghĩ về chuyến đi chơi cuối tuần. Kuon đứng ở góc phòng, khuôn mặt vô hồn, hai tay khoanh trước ngực, cố tỏ ra như đang chú ý - nhưng thực chất cậu đang tính xem còn bao lâu nữa thì tan làm.

``Nghe cũng có lý nhỉ,''  nàng Phù thủy cảm thán, tỏ vẻ hứng khởi với kế hoạch này. ``Sao em không nghĩ ra từ sớm!? Đổ dầu vào, máy móc sẽ chạy lại, đúng như trong phim!''

Korone liếc nhìn cô nàng Phù thủy, trông chẳng tỏ vẻ hứng thú gì, giọng chán nản: ``Sao em cũng phải nghe chỉ dạy của chị ấy vậy chứ? Em chỉ đến đây chào buổi sáng thôi, giờ lại thành đồng phạm.''

Cậu Tổng chỉ có thể bất lực lên tiếng, giọng đầy chua chát: ``Câu đó phải là anh nói mới đúng đó. Anh chỉ là người ngồi làm việc, tự dưng bị kéo vào vụ này, còn bị mắng là không giúp. Chịu hai đứa luôn''

Cậu theo phản xạ tự nhiên liền hỏi một câu tưởng chừng vô hại, nhưng hóa ra lại mở ra một cánh cửa địa ngục: ``Vậy chúng ta sẽ tìm dầu ở đâu? Roboco-chan cần dầu máy loại đặc biệt, không phải loại nào cũng dùng được.''

Không chút do dự, nàng Trung học tươi cười, giọng chắc nịch như một chuyên gia dầu nhớt: ``Cái đó thì không phải lo! Chị nhìn thấy một chai dầu thế này trong nhà bếp đó! Để chị lấy cho!''

Cô chạy biến vào bếp, rồi nhanh chóng trở ra với một chai thủy tinh xanh, bên trong là chất lỏng màu vàng óng, long lanh dưới ánh đèn.

Nàng Phù thủy nghiêng đầu khó hiểu, trong khi cậu Tổng và nàng Khuyển đồng loạt đổ mồ hôi hột, linh cảm chẳng lành.

Shion cầm chai lên, đọc nhãn, giọng đầy hoang mang: ``Đây là\dots?''

Haato cười tươi, giơ ngón tay cái: ``Dầu ô-liu đó! Hàng nhập khẩu, chất lượng cao, chị dùng để trộn salad mỗi ngày. Cũng là dầu mà, chẳng qua là dầu ăn thôi, nhưng chắc cũng bôi trơn được.''

Cả Kuon và Korone đồng loạt tuyệt vọng, hai giọng nói hòa làm một: ``Cái gì đây hả trời\dots''  rồi họ nhìn nhau, cùng thở dài.

\img{5-4h} (giả định)

Nàng Chó nhăn mặt, quay sang cậu Tổng với ánh mắt van nài, đôi tai chó cụp xuống: ``Kuon-senpai, em có thể về được không? Em thấy mình không cần thiết ở đây lắm. Em chỉ muốn về xem phim hoạt hình thôi.''

Tất nhiên, Haato không chịu bỏ cuộc dễ dàng. Cô nhất quyết bắt cả ba người còn lại đứng nhìn mình thực hiện ``phép màu'' này, như một ảo thuật gia trước đám đông.

``Đứng nhìn đây này! Cảm nhận sự kỳ diệu của khoa học và dầu ô liu!''

Cô cẩn thận cầm chai dầu ô-liu mang nhãn hiệu tên mình (có in hình mặt cười), từ từ rót vào\dots\ thứ mà cô cho là bình chứa dầu của tiền bối mình. Thực chất, đó là một khe hở giữa các tấm ốp lưng của Roboco, nơi chất lỏng màu vàng cam vẫn đang rỉ ra.

``Đổ một chút dầu vào đây\dots\ Thế là xong! Roboco-san sẽ tỉnh dậy như trong phim khoa học viễn tưởng!''

\img{5-4i}

\dots\ Ba giây im lặng.

\dots\ Mười giây.

\dots\ Ba mươi giây.

Không có gì xảy ra cả.

Không ánh sáng, không tiếng động, không có màn hình khởi động. Thứ chất lỏng màu vàng cam kia vẫn tiếp tục rỉ ra, chẳng thèm quan tâm đến dầu ô liu mới. Thậm chí nó còn chảy nhanh hơn, như thể đang phản ứng hóa học với loại dầu nhập khẩu.

Nàng Phù thủy hét lên, mặt tái xanh như xác sống, hai tay ôm đầu hoảng loạn: ``Kết thúc rồi! Cả bốn bọn mình sẽ phải vào nhà tù bóc lịch mất thôi! Đổ dầu ăn vào người máy - đó là tội ác chống lại loài người!''

Cô tung chiếc mũ phù thủy của mình lên cao, xoay tít trên không, rồi nó rơi đúng vào mặt cô, che kín tầm nhìn. Cô chới với, suýt ngã.

Nàng Khuyển bàng hoàng khi cô ấy chỉ điểm cả cô vào vòng nguy hiểm: ``Cả tôi cũng có tội à!? Tôi chỉ đứng xem thôi!''

Cậu Tổng cũng bức xúc không kém, giọng vang vọng cả phòng: ``Sao hai đứa lại lôi cả anh vào nữa!? Anh còn chẳng động vào chai dầu! Anh chỉ là khán giả bất đắc dĩ!''

\img{5-4k}

Trong khi cả nhóm đang náo loạn, Haato - kẻ chủ mưu của thảm họa - ung dung định\dots\ chuồn khỏi hiện trường. Cô lén lút quay lưng, nhón chân, rón rén tiến về phía cửa, như một con mèo vừa làm vỡ lọ hoa.

``Thôi, chúc mọi người may mắn, em đi về nhà đây. Hẹn gặp lại vào kiếp sau\dots''

Nhưng nàng ``Tỏi'' không để yên. Shion lập tức lao tới, ôm chầm lấy cả thân người đàn chị, kéo cô lại, quát: ``Không được! Chị cũng phải ở lại với bọn em chứ! Dù có phải lên thiên đường, xuống địa ngục - cả bọn cùng đi!''

\img{5-4l}

Và thế là cuộc chiến đổ tội bắt đầu.

Haato và Shion bắt đầu cãi chem chẻm, vừa tức giận hét vào mặt nhau vừa chỉ điểm, chẳng khác gì hai chính trị gia trong một cuộc tranh luận truyền hình.

Haato phản kháng, giọng cao vút: ``Này, do em gây ra hết đấy! Nếu em không động vào Roboco-san, thì đã chẳng có chuyện gì! Chị chỉ giúp thôi mà!''

Nhưng Shion vẫn không chịu thua, đanh thép đáp lại, tay chỉ thẳng vào mặt đàn chị: ``Ai bảo chị đổ dầu ô-liu vào trong đó!? Có ai đổ dầu ăn vào người máy bao giờ không!? Đổ nước lọc còn khôn hơn!''

Haato lên giọng, cố gắng tự bào chữa, hai tay chống nạnh: ``Này nhé, chính em đã hỏi chị: `có cách gì không?'. Và đó là cách chị làm đấy! Chị đã nghĩ ra, còn các người thì chẳng nghĩ được gì!''

Tất nhiên, lời phân trần đó không thuyết phục được nàng Phù thủy. Shion gào lên: ``Nhưng đó không phải là cách! Chị có hàng tá cách giải quyết khác tốt hơn mà - gọi thợ, gọi cứu hộ, gọi bảo hành - sao lại không dùng!?''

Thấy rằng mình đã đuối lý, Haato nhăm nhe tới nàng Khuyển và cậu Tổng - những người đang đứng ngoài cuộc như khán giả xem phim. Cô chỉ thẳng vào cậu quản lý, giọng đầy ``công lý'': ``Mà\dots\ Thôi, có khi lại do anh ấy gây ra đấy! Anh ấy là quản lý, phải chịu trách nhiệm về mọi thứ trong văn phòng!''

Dần dà, nàng tóc vàng bắt đầu chuyển hướng trách nhiệm sang\dots\ những người không liên quan. Kuon ngỡ ngàng đến mức sốc, miệng há hốc: ``\vie{Ơ cái định công mệnh!?} Sao mấy đứa lại chỉ vào anh!? Anh còn chẳng đả động gì tới chuyện này cả!''

Nàng Trung học thản nhiên đáp lại, theo sau là nàng Phù thủy vội gật gù đồng ý mà không suy xét gì: ``Anh cũng phải có một phần trong đó nữa. Ít nhất thì anh đã không ngăn cản chúng em.''

Câu trả lời của cô càng làm cậu phẫn nộ, cậu bắt đầu buông những ``lời vàng ý ngọc'' về thẳng cả hai cô gái, tay chỉ lung tung: ``\eng{Ơ, đồng cây mát mẻ cái gì đấy!?} Liên quan vãi! Anh có biết các em định làm gì đâu mà ngăn!''

Nàng Phù thủy tiếp tục đổ thêm dầu vào lửa, mặt đầy vẻ ``chân lý'': ``Anh cũng là người cùng với bọn em làm tràn dầu ra còn gì! Khi nãy anh đã nếm thử chất lỏng, chứng tỏ anh có dính líu!''

Cậu trợn mắt nhìn hai người, bất mãn với những lập luận vô lý: ``Này nhá, từ nãy tới giờ anh chỉ đứng đây xem thôi đấy! Nếm thử để xác định chất lỏng chứ có đổ vãi đâu! Mấy đứa hay thật đấy nhở? Anh đây muốn nói với mấy đứa cái này mà có chịu nghe--''

Haato lại tiếp tục cắt ngang lời cậu với một câu vu khống đầy nghệ thuật: ``Bọn em không cần biết! Anh phải là người chịu trách nhiệm! Tìm cách giải quyết đi, nếu không thì lát A-chan vào sẽ tận mắt chứng kiến đấy!''

Tới đây, Kuon bắt đầu tỏ vẻ ngán ngẩm với những lời ``mất não'' từ hai cô gái. Cậu cao giọng, đưa tay lên trán: ``Ơ cái bọn này? Bây giờ không biết đổ cho ai nên lại chuyển sang đổ cho những người không liên quan hả? Hay thật đấy!''

Cậu chỉ tay thẳng vào Korone - người từ nãy đến giờ chỉ đứng nhai kẹo và thỉnh thoảng lắc đầu - và tiếp tục: ``Nếu mà nói như em thì sao em không đổ tội đó cho Koro-chan luôn đi? Em ấy cũng có mặt ở đây từ đầu, chẳng làm gì cả, thế thì cũng đáng bị đổ tội như anh!''

Nàng Khuyển thấy mình bị réo tên lần thứ hai, cô nàng hoảng hốt, suýt rơi cây kẹo: ``Ể!? Sao lại là em!? Em chỉ đến chào buổi sáng thôi! Em không liên quan gì đến vụ dầu dã này cả!''

Cậu hất cằm về phía chiếc giá gỗ với cây thánh giá và đinh mà cô nàng mang tới trước đó: ``Tự dưng em cũng bày ra cái trò đóng đinh đấy còn gì! Đưa thánh giá với đinh ra dọa người ta, bảo thú tội - thế thì cũng có một phần trách nhiệm!''

Nàng Chó sững người, rồi lập tức phản bác, mặt đỏ lên: ``Ơ? Em là người được hai người này giải thích mọi chuyện, xong rồi em mới giúp chứ? Mọi người không thích cách của em thì thôi, sao lại đổ cho em? Em chỉ đùa thôi mà!''

Cô bất mãn, gàm gừ với Haato và Shion: ``Mà em thấy nãy giờ hai chị này kèn cựa với nhau lắm á! Cãi nhau ỏm tỏi, suýt đánh nhau nữa! Các chị tự giải quyết đi, đừng kéo em vào!''

``Này nhá\dots''

Và thế là, cả bốn người tiếp tục tranh cãi, mỗi người một tội, mỗi người một lý do. Ba người con gái thậm chí còn lao vào đánh nhau - kiểu túm tóc, kéo áo, chỉ mặt nhau - khiến cậu Tổng phải vất vả lắm mới kéo Haato, Shion và Korone ra khỏi màn ``giao lưu võ thuật'' bất đắc dĩ. Cậu vừa giữ những cái đầu nóng của họ, vừa giữ bình tĩnh cho bản thân.

Cuối cùng, khi mọi người đã mệt lử, thở hồng hộc, cậu chỉ còn biết thở dài trong bất lực, ngả người ra ghế, ngân nga một câu hát nhẹ nhàng, đầy triết lý: ``\vie{Ôi\dots\ Đàn bà là những niềm đau\dots}''

\ct

Trong lúc bốn người họ còn đang tranh cãi kịch liệt, thậm chí cả nàng Chó và cậu Tổng cũng bị lôi vào cuộc chiến (với những lý do vô cùng vô lý), thì cánh cửa bất ngờ mở ra.

Một bóng dáng quen thuộc bước vào. Roboco  xuất hiện với vẻ mặt hớn hở, hai tay ôm một túi đồ, mắt sáng rực, vô cùng tươi tỉnh. Trông cô chẳng có vẻ gì là vừa bị rỉ dầu hay hết pin cả. Cô cất giọng chào, vang dội khắp phòng: ``Mình về rồi đây! À-rế? Sao mọi người đông vui thế? Có tiệc à?''

\img{5-4m}

Cả bốn người đồng loạt đứng sững lại. Haato đang túm cổ Shion, Korone đang cắn tay Kuon (khi cậu cố gắng kéo cô ra), tất cả đều quay đầu nhìn về phía cánh cửa. Và rồi, họ chứng kiến một cảnh tượng không thể tin nổi:

Roboco - bằng xương bằng thịt, chính xác hơn là bằng kim loại - đang đứng trước mặt họ, tươi cười, khỏe mạnh, và hoàn toàn\dots\ không bị rỉ dầu.

Trước mắt cô là cảnh tượng hỗn loạn: ba cô gái mặt mày xây xát, quần áo xộc xệch, tóc tai rối bù như vừa trải qua một trận chiến sống còn trong rừng rậm. Còn cậu con trai thì thở dốc, bám lấy mép bàn để đứng vững, tay áo bị cắn rách một mảng.

Nhìn thấy Roboco, cậu Tổng như bắt được cứu tinh sau cơn bão, hai mắt rưng rưng, giọng nghẹn ngào: ``May quá em về rồi\dots\ Giải quyết giùm anh đống này cái! Chỉ cho họ biết cái thứ đang nằm trên sàn kia là gì đi!''

Lúc này, cả hai ``kẻ phạm tội''  Haato và Shion - sững sờ nhìn Roboco thật. Cô nàng người máy bằng xương bằng thịt (chính xác hơn là bằng kim loại) đang đứng trước mặt họ, khỏe mạnh, tươi tỉnh. Còn Korone thì nhanh trí chạy ngay về phía cậu quản lý, tránh xa khỏi hai đồng phạm, như thể tuyên bố: ``Tôi không liên quan, tôi chỉ là khán giả!''

Haato lắp bắp, tay run run chỉ về phía Roboco: ``À-rế\dots? Nhưng nếu Roboco-san đang ở đây, thì\dots''

Shion cũng hoảng hốt, tiếp lời: ``\dots\ cái thứ nãy giờ bọn mình lau chùi, cãi cọ, đổ dầu ô liu vào\dots\ là cái gì vậy?''

Cả hai quay ngoắt sang nhìn nhau, mặt tái xanh.

Roboco chớp mắt vài cái, nhìn quanh, rồi cuối cùng cũng hiểu ra vấn đề. Cô tiến lại gần ``xác'' đang nằm trên sàn, nhìn nó một lượt, rồi bật cười. Vẫn giữ giọng hí hửng như thường lệ, cô vui vẻ giải thích, giọng đầy tự hào: ``Ồ, mấy đứa nghĩ sao thế? Đó là bình chứa cà ri mới của chị đó! Nó có thể đựng được nhiều cà ri lắm\~{} Chị đặt làm riêng, chất liệu cao cấp, chống tràn, chống sốc, lại còn có hình dáng giống người để dễ mang theo. Hôm qua chị dùng nó để đựng cà ri tắm, xong định mang về nhà rửa, ai ngờ bỏ quên ở đây.''

Cả căn phòng bỗng trở nên yên lặng đến mức có thể nghe thấy tiếng ve kêu ngoài cửa sổ.

Haato và Shion hoàn toàn ngớ người, mắt mở to, miệng há hốc. Họ quay sang nhìn nhau, rồi nhìn xuống thứ chất lỏng màu vàng cam đang chảy loang trên sàn - thứ mà họ đã lau chùi, đổ dầu ô liu vào, tranh cãi đổ tội, suýt đánh nhau, và bây giờ thì biết đó chỉ là\dots\ cà ri từ một cái bình chứa hình người.

Cậu quản lý thở dài bất lực, giơ tay lên trời: ``Đấy, cái anh định nói với mấy đứa từ lúc đầu là như vậy đấy! Anh bảo rồi, đó là cà ri, anh đã nếm thử rồi, thế mà có ai nghe đâu!''

Sau đó, cậu quay sang Roboco - người lúc này đang cúi xuống nhìn cái bình chứa tội nghiệp với vẻ mặt tiếc nuối - và bồi thêm một câu: ``Với cả Robo-chan ạ, cái bình mới của em sao trông nó dị thế? Giống hệt người thật, còn có tóc, có mắt, lúc nằm im trông y như xác chết. Em đặt ở đâu vậy?''

Roboco cười hì hì: ``Em đặt trên mạng đó, shop chuyên làm bình đựng đồ cosplay. Họ bảo là hàng độc, không giống ai, em thấy cũng ổn mà.''

Bầu không khí trở nên căng thẳng một cách khó tả. Haato và Shion từ từ quay sang nhìn nhau, ánh mắt đanh lại, đôi môi mím chặt. Một tia chớp lóe lên trong mắt họ.

``\textbf{Lừa cả bọn này rồi!}''  cả hai đồng loạt hét lên, giọng đầy phẫn nộ như thể vừa bị lừa bán một món hàng giả.

Không một chút do dự, họ lao tới cái bình chứa xấu số, cùng lúc nâng nó lên. Trong tích tắc, với tất cả sức lực còn lại sau một buổi sáng đầy biến động, họ ném thẳng nó ra ngoài cửa sổ.

\img{5-4o}

Kính vỡ tan tành, những mảnh vỡ lấp lánh rơi xuống đường dưới ánh nắng. Cái bình cà ri hình người bay vút ra ngoài, vẽ một đường cong hoàn hảo trong không khí, rồi biến mất sau những tán cây.

Roboco hét lên, giọng đầy thất vọng, hai tay ôm đầu: ``\textbf{Bình chứa của chị!!} Cái bình mà chị đã đặt hàng từ ba tháng mới có, giờ bay mất rồi!''

Cô chạy ra cửa sổ, nhìn theo, chỉ kịp thấy chiếc bình đáp xuống thùng rác bên dưới một cách\dots\ hoàn hảo.

Haato và Shion thở hổn hển, đứng nhìn ô cửa sổ bị phá tan tành. Họ dường như vừa mới nhận ra mình vừa làm gì.

Cậu quản lý chỉ biết nhìn ô cửa sổ mà cười một cách cay đắng, lắc đầu: ``Một sự hiểu lầm cực nặng ở đây\dots\ Bình chứa cà ri bị tưởng là Roboco chết, đổ dầu ô liu vào rồi đem ném ra ngoài.''

Korone từ phía sau mon men tới, nhìn mọi người, rồi thận trọng lên tiếng: ``Thế\dots\ bây giờ mình có cần dọn kính vỡ không? Hay là gọi thợ? Hay là\dots\ chơi bài chuồn đây?''

Roboco quay lại, mặt vẫn hơi buồn, nhưng cô thở dài, đưa tay quệt quệt: ``Thôi, cũng tại chị để bình ở đấy. Lần sau chị sẽ để trong tủ. Còn về cửa sổ\dots\ chúng ta cùng dọn đi, kẻo lát có người vào lại thắc mắc.''

Haato và Shion cúi đầu, lí nhí: ``Em xin lỗi\dots\ Lần sau tụi em sẽ hỏi kỹ trước khi hành động.''

Kuon thở dài lần cuối cùng trong ngày, vỗ vai Roboco: ``Thôi, dọn đi. Rồi gọi thợ sửa cửa sổ. Tự dưng một buổi sáng đẹp trời lại biến thành trận chiến vì bình cà ri.''

Cả năm người bắt tay vào dọn dẹp. Haato và Shion lau sàn, Korone nhặt mảnh kính, Roboco lau người cho cái bình còn sót lại (đã vỡ một góc), còn Kuon thì quét rác. Ngoài cửa sổ, gió thổi nhẹ, ánh nắng tràn vào, và từ xa vọng lại tiếng ve kêu râm ran - như chẳng có gì xảy ra.
%==========================================TRUYỆN PHỤ=========================================
\omk

Roboco không mất nhiều thời gian để lấy lại cái bình - cô chạy như bay xuống tầng, mò vào thùng rác, và lôi nó lên trước sự chứng kiến của mấy người đi đường. Nhưng khi cô quay trở lại văn phòng, khuôn mặt đã không còn vui vẻ như trước. Cái bình vẫn còn sử dụng được, nhưng giờ nó đã có vài vết xước do va đập. Trông nó như một chiến binh vừa trở về từ chiến trường, đầy thương tích.

``Mồ\dots\ Cái bình này đắt tiền lắm ớ\dots\ Chị đã để dành cả tháng lương part-time đấy\dots''  Roboco than thở, vỗ vỗ lên chiếc bình như vỗ về một đứa trẻ bị thương.

Trong khi đó, cậu Tổng - giờ đang bôi thuốc sát trùng cho ba kẻ vừa trải qua trận chiến không khoan nhượng - chỉ biết thở dài ngao ngán. Trên tay cậu là một lọ cồn và bông gòn, trước mặt là Haato, Shion và Korone đang ngồi xếp hàng, mặt mày xây xát, tay chân trầy xước.

Roboco nhìn cảnh tượng, tò mò hỏi: ``Thế, có chuyện gì mà ba người họ xây xát mặt mày vậy, Kuon-senpai? Nhìn như vừa đánh nhau với gấu ý.''

Cậu chậm rãi kể lại toàn bộ câu chuyện dở khóc dở cười cho Roboco nghe, từ lúc ``dầu'' tràn ra (thực ra là cà ri), cho đến khi Haato quyết định đổ dầu ô-liu vào (với lý do ``cũng là dầu mà''), và màn hiểu lầm đỉnh cao nhưng ngớ ngẩn cuối cùng: cả nhóm cãi nhau, đổ tội, suýt đánh nhau, chỉ vì một cái bình giống người.

Roboco khoanh tay, gật gù như một vị thẩm phán sau khi nghe xong lời khai: ``Em hiểu rồi\dots\ Mọi người tranh cãi nhau chỉ vì cái bình này nhìn giống em chứ gì? Tóm lại là do chị quên bình ở đây, các em thấy giống người thật, tưởng chị bị gì, rồi hoảng loạn?''

Cả bốn người lập tức gật đầu như máy. Cậu Tổng liếc về phía Haato và Shion, nói trong sự ngán ngẩm: ``Giá như mà em nói trước cho họ thì mọi chuyện đâu có như thế này. Chỉ cần một câu `đó là bình đựng cà ri của chị' là xong, đằng này\dots''

Nàng Song tính và nàng Phù thủy cũng gật đầu lia lịa, mặt đầy vẻ ``đúng đúng, tại bà hết''.

``Đúng đó, cái bình của bà nhìn dễ hiểu nhầm quá! Nằm im trên sàn, ôm mặt, còn có tóc - trông y hệt xác chết! Ai chẳng hoảng!''  Haato buông một câu xanh rờn.

``Bọn em còn tưởng đó là chị đang nằm bất tỉnh, hết pin, rỉ dầu các thứ. Nếu biết là bình, tụi em đã chẳng lau chùi, đổ dầu, rồi cãi nhau ỏm tỏi.''  Shion tiếp lời.

Roboco cười trừ, gãi đầu đầy áy náy: ``Rồi, rồi, chị xin lỗi\dots\ Lần sau chị sẽ dán nhãn `bình đựng cà ri' lên trán nó, khỏi ai nhầm.''

Cậu Tổng vẫn chưa từ bỏ thắc mắc ban đầu của mình. Cậu vừa bôi cồn lên vết xước trên tay Haato vừa hỏi: ``Mà em vẫn chưa trả lời câu hỏi của anh đâu. Cái bình đó ở đâu ra vậy? Trông nó như hàng độc, không phải mua ở tiệm thường.''

Roboco lập tức sáng mắt lên, giọng đầy tự hào, như một nhà sưu tập vừa tậu được bảo vật: ``À, cái đó em săn được ở chợ ảo Akihabara đấy! Shop chuyên bán đồ cosplay, họ có mấy cái bình hình người rất độc. Anh thấy sao? Hợp với em không?''

Cậu nhìn cái bình - thứ vừa gây ra một chuỗi sự kiện hỗn loạn từ lúc Haato và Shion tới đây đến giờ - rồi thở dài phán: ``Trông nó cũng dị đấy. Mà em mua bởi vì nó giống em đúng chứ? Giống y hệt, chỉ khác là nó không nói được.''

Roboco ``ưm!'' một tiếng đầy đắc ý, hai tay ôm bình vào lòng, mắt sáng rực. Nhưng rồi cậu Tổng chỉ lắc đầu, thêm một câu chốt hạ: ``Đúng chuẩn PON luôn. Mua cái bình giống mình, để rồi người khác tưởng mình chết, suýt đánh nhau. PON không thể chữa.''

Lập tức, cô nàng bật dậy phản đối, mặt đỏ lên, phồng má giận dỗi, trông chẳng khác gì một đứa trẻ bị chọc quê: ``\textbf{Em không có PON mà!} Em chỉ vô tình thôi! Chỉ có một lần này thôi!''

Cả phòng phá lên cười. Ngay cả Haato và Shion, dù đang bị bôi thuốc đau điếng, cũng không nhịn được.

Sau khi mọi chuyện đã lắng xuống, cậu Tổng khẽ hắng giọng, đứng dậy, vẻ mặt trở nên nghiêm nghị như một kế toán đang tổng kết sổ sách cuối năm: ``Bây giờ anh bắt đầu tính thiệt hại trong ngày hôm nay nha. Các em ngồi yên lắng nghe.''

Haato và Shion nghe vậy, liền phản ứng lại, giọng đầy nghi ngờ: ``Cái gì nữa vậy!? Chẳng lẽ còn phạt nữa?''

Kuon không thèm giải thích, cậu tiếp tục dõng dạc tuyên bố, giống như một vị quan tòa tuyên án: ``Tiền cửa sổ, mỗi người phạt 1500 yên. Kính vỡ bay ra ngoài đường, may không trúng ai, nhưng phải thay mới.''

Cả hai nuốt nước bọt, không ai dám lên tiếng phản đối. 1500 yên - cũng bằng một bữa ăn trưa, nhưng vẫn xót.

``Tiếp theo, tiền tổn thương thể xác và tinh thần, mỗi người trừ 20\% lương!''

Ngay lập tức, cả hai người vùng lên như lò xo, mặt đỏ bừng, mắt trợn ngược: ``\textbf{Gì mà tiền tổn thương thể xác với tinh thần!?} Bọn em có làm gì đâu! Chỉ là hiểu lầm thôi mà!''

Cậu lại hắng giọng, nhìn họ với ánh mắt đáng sợ - thứ ánh mắt mà chỉ một lần là đủ để bất kỳ nhân viên nào cũng phải im bặt: ``Đã nói hết đâu mà gân cổ lên? Im lặng, nghe tiếp.''

Hai cô gái vội ngồi xuống, hai tay đặt lên đầu gối, im như tờ.

Cậu trở lại giọng bình thường, nhưng vẫn đầy uy quyền: ``Do hai người đã đánh nhau - anh có camera ghi lại hết - nên số tiền này coi như bù trừ 10\% cho nhau. Kết quả cuối cùng: mỗi người chỉ phải nộp 10\% tiền lương tháng này.''

Cả hai thở phào nhẹ nhõm, nhưng vẫn cảm thấy không cam tâm. Shion băn khoăn, rụt rè hỏi: ``Mà, 10\% đó dùng vào việc gì? Không phải tiền phạt vớ vẩn chứ?''

Cậu đáp, giọng tỉnh bơ nhưng đầy lý lẽ: ``5\% là tiền thuốc men - bôi tróc da, băng vết thương, mua cồn, bông gòn. 5\% còn lại là phí tổn thất tinh thần cho Koro-chan - người bị lôi vào vụ này một cách oan uổng, bị mắng, bị chỉ trích, suýt bị đóng đinh.''

Korone nghe thấy tên mình, bỗng dưng mỉm cười, vẫy đuôi vui vẻ. Cô gật gù ra vẻ ``đúng đúng, tôi bị ảnh hưởng nặng nề lắm đấy''.

Haato lại nhảy dựng lên, định phản đối: ``Cái gì chứ!? Này nhé - em có làm gì Korone-chan đâu, tự nhiên trừ lương của em để đền cho cô ấy--''

``Cấm cãi,''  Kuon cắt ngang, giọng đầy uy quyền, mắt nheo lại. ``Cãi nữa trừ thêm lương. Anh quá mệt với hai đứa rồi. Sáng giờ chỉ lo dọn dẹp, bôi thuốc, làm trọng tài. Anh còn chưa tính công anh đấy.''

Nhìn khuôn mặt đáng sợ khi tức giận của cậu, nàng Trung học chỉ có thể ỉu xìu xuống nước, lí nhí: ``V-Vâng\dots\ Bọn em biết lỗi rồi\dots''

Sau khi tính xong tiền phạt, Kuon trông mệt mỏi hẳn, hai mắt thâm quầng như vừa thức trắng đêm. Hai ``phạm nhân''  Haato và Shion - thì tiu nghỉu, mặt cúi gằm, lòng thầm tiếc số tiền lương sắp bị trừ. Còn nàng Khuyển Korone thì đã bình tĩnh trở lại, vẫy đuôi vui vẻ, vớicảm giác như vừa được đền bù xứng đáng.

Roboco nhìn cảnh tượng, khẽ cười, rồi đề xuất: ``Dù gì mấy đứa cũng đã mệt rồi, sao không thử cà ri mà chị làm nhỉ? Chị nấu từ hôm qua, hôm nay định mang đến cho cả nhóm.''

Cả bốn người đồng thanh, giọng đầy hào hứng - sự hào hứng duy nhất trong cả buổi sáng:

Haato: ``Được chứ!''
Shion: ``Vâng\~''
Korone: ``OK!''
Kuon: ``Đâu, đâu? Có gì ăn là anh xốc ngay.''

Roboco cười tủm tỉm, rồi nói tiếp: ``Dạng uống!''

Cả bốn người sững người. ``Cà ri dạng uống?''  Kuon nhướng mày.

Roboco hào hứng giải thích, cô vừa nói vừa lấy mấy cái cốc nhựa ra: ``Mấy đứa cứ cầm cốc ra, để chị chỉ. Cà ri này được nấu đặc biệt, xay nhuyễn, nấu với nước dùng, uống vào mát ruột, bổ phế, lại giải nhiệt mùa hè.''

Cô quay ra sau ``lưng'' của cái bình - nơi có một vòi nhỏ - và lần lượt rót cho mỗi người một cốc cà ri nóng hổi, màu vàng cam bắt mắt, tỏa hương thơm nức mũi của nghệ, thì là, và các loại gia vị.

\img{5-4p}

Khi cả bốn người dùng thử, ai nấy đều tấm tắc khen ngon. Kuon nhấp một ngụm, mắt sáng lên: ``Ừm, ngon thật! Uống cũng lạ, nhưng mà ngon. Robo-chan, em làm được đấy.''

Korone thì bốc hết cốc một hơi: ``Thêm cốc nữa!''

Haato và Shion cũng xuýt xoa không ngớt. Nhưng rồi, khi cốc cà ri đã vơi đi một nửa, Haato bỗng nhớ ra một điều\dots\ khủng khiếp.

Cô quay sang Shion, ghé sát tai, thì thầm với giọng đầy lo lắng: ``\hl{\dots\ Có nên nói cho chị ấy không nhỉ? Vụ dầu ô liu ấy\dots\ Chị ấy đang uống cà ri từ cái bình mà mình đã đổ dầu vào.}''

Shion mắt mở to, rồi kịp thời lấy lại bình tĩnh. Cô cũng thì thầm đáp, giọng như mật: ``\hl{Tốt nhất là đừng có nói. Im lặng là vàng.}''

Haato gật đầu, nuốt nước bọt, rồi cố tỏ ra tự nhiên, nâng cốc lên uống tiếp. ``Không sao đâu, dầu ô liu cũng tốt cho sức khỏe mà\dots\ phải vậy không?''

Shion gật gù, nhưng mặt thì hơi xanh.

Roboco vẫn rót cà ri, vẫn cười hồn nhiên, không hề hay biết. Cô đưa cốc lên miệng, vừa uống vừa khen: ``Cà ri của chị ngon nhỉ? Bí quyết là cho thêm sữa tươi và một chút mật ong.''

Cả Haato và Shion nhìn nhau, cố nín cười, mà lòng thì thầm cầu nguyện: ``Mong chị ấy không phát hiện ra vị dầu ô liu\dots''

Mọi chuyện cuối cùng cũng đã trở về yên bình. Ngoài cửa sổ, nắng vàng trải dài, ve vẫn kêu râm ran. Roboco vui vẻ chia cà ri, Kuon thảnh thơi uống, còn Haato, Shion, Korone thì đã bỏ qua những hiểu lầm, ngồi quây quần bên nhau, tiếng cười lại vang lên trong căn phòng.

\ct

Sau một hồi cười đùa và những ngụm cà ri cuối cùng, không khí trong phòng đã trở nên nhẹ nhàng hơn hẳn. Cậu Tổng ngồi tựa vào ghế, nhìn quanh, rồi bất chợt quay sang Korone với một câu hỏi mà từ nãy giờ vẫn ám ảnh cậu: ``Mà này, Koro-chan.''

``Gì vậy?''  nàng Chó ngước lên, miệng vẫn còn đang nhai kẹo.

``Tưởng mũi em thính lắm mà? Sao em không nhận ra cái mùi cà ri này từ đầu? Nếu em phát hiện sớm, thì đâu có chuyện hỗn loạn đến thế.''

Korone cười khì, đưa tay chỉ mũi: ``Em bị tịt mũi ạ\dots\ Hôm qua ngủ điều hòa hơi nhiều, sáng dậy thấy nghẹt mũi, chẳng ngửi thấy gì hết.''

Cậu chỉ biết thở dài lần cuối trong ngày hôm nay, đưa tay ôm trán, chán chẳng buồn nói: ``Thật luôn hả giời\dots\ Bị tịt mũi vào đúng ngày quan trọng nhất, khi cả văn phòng sắp bốc cháy vì cà ri. Số nhọ vãi.''

Sau câu nói của cô, cả ba người còn lại - Haato, Shion và Roboco - chỉ biết nhìn nhau trong sự ngơ ngác. Cậu con trai vừa thở dài vừa lẩm bẩm: ``Bảo sao từ nãy giờ em cứ bình chân như vại, không hề tỏ ra khó chịu, mặc cho mùi cà ri nồng nặc. Hóa ra là chẳng ngửi thấy.''

Haato và Shion cười trừ, cùng lên tiếng: ``Ừ nhỉ\dots\ Tại sao bọn này không nghĩ ra nhỉ? Đơn giản thế mà.''

``Cái mùi cà ri đặc trưng thế cơ mà, nếu không bị tịt mũi thì ai chẳng phát hiện ra.''

Korone lại nhấp một ngụm cà ri, đưa cốc lên môi, rồi bỗng nhiên mắt cô mở to, lông mày nhướng lên: ``Khoan\dots\ hình như em bắt đầu ngửi được mùi rồi!?''

Cả ba người lập tức giật mình nhìn về phía nàng Khuyển, trong khi cô nàng dường như đang định hình lại hương vị của thứ mình vừa uống. Không khí trở nên im lặng trong vài giây - một sự im lặng nặng nề, đầy kịch tính.

Rồi Korone hét lên, giọng xé toạc bầu không khí: ``\textbf{Ối giời ơi!? Sao vị nó lạ thế!? Có gì đó không ổn lắm!}''

Lập tức, cô đặt cốc xuống, chạy như bay ra bàn nước, tay với lấy bình lọc nước, mở vòi uống ừng ực. Cả bọn chỉ biết ôm bụng cười ngặt nghẽo, nước mắt chảy ra vì cười nhiều hơn là vì\dots\ cà ri.

Kuon đứng dậy, vỗ tay hai cái thật lớn: ``Rồi rồi, kết thúc rồi! Giờ ai về nhà nấy nghỉ ngơi đi. Hôm nay thế là đủ rồi! Anh không muốn thêm bất kỳ một hiểu lầm nào nữa.''

Haato, Shion, và Roboco đều gật đầu đồng ý, dù trong lòng vẫn còn luyến tiếc chút ít về chuyện vừa xảy ra. Nhưng ít nhất, mọi chuyện cũng đã ngã ngũ. Không ai bị thương nặng, không ai mất mạng, và cái bình cà ri vẫn còn nguyên - chỉ thêm vài vết xước làm kỷ niệm.

Trước khi rời đi, Roboco vẫn không quên nhìn lại cái bình chứa yêu quý của mình. Cô vuốt nhẹ lên nó như thể đang dặn dò một đứa trẻ: ``Lần sau chị sẽ để em ở nơi dễ nhận biết hơn\dots\ Có thể dán ảnh chị lên người em, hoặc viết chữ `bình đựng cà ri' to đùng. Kẻo lại có đứa tưởng nhầm mất\~{}''

Cặp đôi phá phách - Haato và Shion - cùng quay sang lườm cô một cái, nhưng chẳng ai nói gì nữa. Họ chỉ biết cười trừ, rồi nhanh chóng thu dọn đồ đạc.

Và thế là, ngày hỗn loạn với vô số tình huống oái oăm đã chính thức khép lại, để lại cho cả nhóm một bài học đáng nhớ: **không nên tự ý suy diễn lung tung khi chưa kiểm tra kỹ**, và cũng là một kinh nghiệm quý báu trong việc **phân biệt bình chứa với người máy sống** - một kỹ năng mà có lẽ chẳng ai ngờ mình cần.

\ct

Sau khi mọi người đã rời đi, văn phòng trở nên yên tĩnh. Cậu Tổng thở dài một hơi, vừa kiểm tra lại danh sách công việc trên điện thoại vừa lầm bầm: ``Rồi, xong vụ phân biệt bình chứa với người máy sống\dots\ Hy vọng tuần sau sẽ bình yên hơn. Nhưng mà tiếp theo thì--''

Cậu chưa kịp dứt lời thì điện thoại rung lên, báo hiệu một tin nhắn mới. Kuon mở ra, thấy một dòng chữ ngắn gọn nhưng đầy khí thế, được viết bằng chữ rất to và gạch chân:

\begin{center}
  ``\underline{\textbf{Này! Văn phòng có khách nha! Chuẩn bị tinh thần đi!!}}''
\end{center}

Không có tên người gửi, không có lời chào, chỉ có một câu cảnh báo như thể một cơn bão sắp ập đến.

Cậu con trai nhướng mày, cảm giác như có một luồng điện chạy dọc sống lưng. Cậu đã trải qua đủ mọi chuyện từ lúc vào làm ở đây, nhưng chưa bao giờ ai nhắn tin kiểu ``chuẩn bị tinh thần'' như thế này.

``\textit{Khách? Khách nào mà phải chuẩn bị tinh thần? Hay là\dots\ khách đặc biệt?}''

\il

Ở đâu đó trong tòa nhà, những tiếng bước chân vang lên dồn dập, xen lẫn tiếng trò chuyện rộn ràng. Không phải một người, mà là cả một nhóm.

Một giọng nói thở dài, đầy vẻ mệt mỏi như đã quá quen với cảnh hỗn loạn: ``Lại đến văn phòng của Kuon-kun nữa\dots\ Lần này nhớ đừng làm gì quá đáng đấy nhé.''

Một giọng khác cười khúc khích, đầy hứng thú, như thể sắp được tham gia một lễ hội: ``Yay! Em muốn gặp anh quản lý huyền thoại đó! Nghe nói ảnh xử lý đủ mọi chuyện điên rồ!''

Một kẻ nghịch ngợm hát một giai điệu ngẫu hứng, chẳng rõ là bài hát gì, nhưng nghe rất vui tai: ``\textit{Lalalala~ hôm nay trời đẹp, đi gặp quản lý, chọc ảnh một trận, rồi về ngủ khò\dots}''

Và một người nữa, với giọng trầm ấm, khe khẽ ngân nga như đang chuẩn bị cho một buổi diễn quan trọng: ``\textit{Hmm\dots\ hmm\dots\ luyện giọng một chút, lát gặp anh ấy chắc sẽ cần hát\dots}''

Màn kịch kế tiếp đã sẵn sàng. Tâm điểm của nó - một vị khách với thân hình to con, tính cách tràn đầy năng lượng, phong thái tự tin và có phần\dots\ ``hơi quá'' - đang sải bước về phía văn phòng. Cô ta không chỉ đến một mình, mà kéo theo sau là **bốn vị khách nữa cùng thế hệ**.

Họ là những người mà Kuon chưa từng gặp mặt trực tiếp, nhưng tiếng tăm của họ đã vang xa khắp công ty: Thế hệ thứ Tư của \hll.

Và họ sẽ trở thành những người tiếp theo mà Hikawa Kuon của chúng ta phải đối đầu, phải quản lý, phải dọn dẹp rác rưởi sau lưng, và có lẽ\dots\ phải học cách sống chung với những cơn bão mang tên ``thần tượng ảo'' trong những khoảng thời gian kế tiếp.

\end{document}